% Generated draft: requires mathematical and source review.
\documentclass[11pt]{article}
\usepackage[margin=25mm]{geometry}
\usepackage{fontspec}
\setmainfont{Latin Modern Roman}
\usepackage{amsmath,amssymb,mathtools,unicode-math}
\setmathfont{Latin Modern Math}
\usepackage{enumitem,needspace,xurl,hyperref}
\setlength{\parindent}{0pt}
\setlength{\parskip}{5pt}
\setlength{\emergencystretch}{3em}
\newcommand{\N}{\mathbb{N}}
\newcommand{\Z}{\mathbb{Z}}
\newcommand{\Q}{\mathbb{Q}}
\newcommand{\R}{\mathbb{R}}
\newcommand{\C}{\mathbb{C}}
\newcommand{\F}{\mathbb{F}}
\newcommand{\A}{\mathbb{A}}
\newcommand{\PP}{\mathbb{P}}
\newcommand{\E}{\mathbb{E}}
\newcommand{\eps}{\varepsilon}
\newcommand{\dd}{\,\mathrm{d}}
\begin{document}
\section*{UnsolvedMath ID: 10400151}
\textbf{Source problem code:} catalogue entry
\subsection{Definitions and mathematical statement}
For \(n\ge2\), the group
\[
 \operatorname{Spin}^c(n)=
 (\operatorname{Spin}(n)\times\mathrm U(1))/\langle(-1,-1)\rangle
\]
maps naturally to \(\mathrm{SO}(n)\). A \(\operatorname{Spin}^c\) structure on an oriented manifold is a lift of its oriented orthonormal frame bundle through this map, with the corresponding stable-tangent and boundary conventions when used in a cobordism category.

Problem 8.5 asks to ``formulate and find \(\operatorname{Spin}^c\) TQFTs.'' A usual formal interpretation, after selecting dimension \(n\) and coefficient field \(F\), would define a \(\operatorname{Spin}^c\) cobordism category with closed decorated \((n-1)\)-manifolds as objects and compatible decorated \(n\)-cobordisms as morphisms. It would then construct symmetric monoidal functors
\[
 Z:\operatorname{Cob}^{\operatorname{Spin}^c}_n
       \longrightarrow\operatorname{Vect}_F,
\]
assigning state spaces and cobordism maps that respect identities, gluing and disjoint union. This is explanatory rather than a source-fixed axiom system: the record does not specify dimension, target category, finite-dimensionality, linearity conventions, anomaly treatment, or the exact handling of boundary structures.

Both the formulation and the construction are requested. The source does not prescribe which \(\operatorname{Spin}^c\) invariants should be recovered, whether nontriviality or detection of the \(\operatorname{Spin}^c\) structure is required, or what distinguishes an acceptable example from a theory that forgets the extra structure. Those criteria must be recovered from context rather than answered by an arbitrarily chosen trivial theory.
\subsection{Short English statement}
Specify a suitable notion of field theory for manifolds equipped with \(\operatorname{Spin}^c\) structures and construct examples. The source gives no dimension or target axioms and no required invariant or nontriviality criterion. Its supplied classification is open-triage.
\subsection{Sources}
\begin{itemize}
\item Local UnsolvedMath registry, permanent ID \(10400151\), exact source code \texttt{catalogue entry}, Problem 8.5. The extracted ``spin c'' is read as \(\operatorname{Spin}^c\). No individual proposer or proposed year is supplied.
\item Ohtsuki, \emph{Problems on invariants of knots and 3-manifolds} (2002), Problem 8.5, PDF page 124: \url{https://msp.org/gtm/2002/04/gtm-2002-04-024s.pdf}. Bold printed headings distinguish 231 numbered items from surrounding references.
\item Extraction pdf-font-structure, accessed 2026-07-29. Source evidence NEEDS\_REVIEW presents an open problem and requires a current primary-literature audit. Supplied research classification OPEN-TRIAGE repeats the formulation-and-construction request without further references.
\item Dataset attribution: Ulam AI / UnsolvedMath Contributors, CC BY 4.0. Source rights NEEDS\_REVIEW: publicly accessible author/publisher copy, redistribution terms requiring review.
\end{itemize}
\end{document}
